\documentclass[11pt]{article}
\usepackage[utf8]{inputenc}
\usepackage[T1]{fontenc}
\usepackage[english]{babel}
\usepackage{amsmath,amssymb}
\usepackage[a4paper,margin=1in]{geometry}
\usepackage{fancyhdr}
\usepackage{hyperref}
\usepackage[protrusion=true,expansion=false]{microtype}
\usepackage{array}

\pagestyle{fancy}
\fancyhf{}
\fancyhead[C]{\small Cayley, Collected Mathematical Papers --- Vol.~III, pp.~129--153 (Papers 188--193)}
\fancyfoot[C]{\thepage}
\renewcommand{\headrulewidth}{0.4pt}

\providecommand{\dd}{d}

\hypersetup{
  pdftitle={Cayley Collected Mathematical Papers, Vol III, pp 129--153},
  pdfauthor={Arthur Cayley (transcription)},
  hidelinks
}

\setcounter{page}{129}

\allowdisplaybreaks

\begin{document}

%===============================================================
% PDF p.151 = book p.129 -- Paper 188 begins
%===============================================================
\begin{center}
\small 188]\hfill\textsc{}\hfill\textsc{129}
\end{center}

\bigskip

\begin{center}
{\Large\bfseries 188.}\\[8pt]
{\large ON THE SIMULTANEOUS TRANSFORMATION OF TWO HOMO-\\
GENEOUS FUNCTIONS OF THE SECOND ORDER.}\\[6pt]
\small [From the \emph{Quarterly Mathematical Journal}, vol.~\textsc{ii}.\ (1858), pp.~192--195.]
\end{center}

\bigskip

\noindent \textsc{In} a former paper with this title, \emph{Cambridge and Dublin Math.\ Journal}, t.~\textsc{iv}.\ [1849], pp.~47--50 [74], I gave (founded on the methods of Jacobi and Prof.\ Boole) a simple solution of the problem, but the solution may I think be presented in an improved form as follows, where as before I consider for greater convenience the case of three variables only.

Suppose that by the linear transformation\footnote{I represent in this manner the system of equations
\[ x = \alpha x_{1} + \beta y_{1} + \gamma z_{1}, \quad \&\text{c.} \]
and so in all like cases.}
\[
(x,\ y,\ z) = \left(\begin{array}{ccc} \alpha, & \beta, & \gamma \\ \alpha', & \beta', & \gamma' \\ \alpha'', & \beta'', & \gamma'' \end{array}\right)\!\!\genfrac{(}{)}{0pt}{}{x_{1},\ y_{1},\ z_{1}}{}
\]
we have identically
\[
(a,\ b,\ c,\ f,\ g,\ h\,\mid\, x,\ y,\ z)^{2} = (a_{1},\ b_{1},\ c_{1},\ f_{1},\ g_{1},\ h_{1}\,\mid\, x_{1},\ y_{1},\ z_{1})^{2},
\]
\[
(A,\ B,\ C,\ F,\ G,\ H\,\mid\, x,\ y,\ z)^{2} = (A_{1},\ B_{1},\ C_{1},\ F_{1},\ G_{1},\ H_{1}\,\mid\, x_{1},\ y_{1},\ z_{1})^{2},
\]
and write also
\[
(\xi,\ \eta,\ \zeta) = \left(\begin{array}{ccc} \alpha, & \alpha', & \alpha'' \\ \beta, & \beta', & \beta'' \\ \gamma, & \gamma', & \gamma'' \end{array}\right)\!\!\genfrac{(}{)}{0pt}{}{\xi_{1},\ \eta_{1},\ \zeta_{1}}{}.
\]

\newpage

%===============================================================
% PDF p.152 = book p.130
%===============================================================
\begin{center}
\small 130\hfill\textsc{on the simultaneous transformation of two}\hfill[188
\end{center}

\medskip

\noindent Comparing these with the relations between $(x,\ y,\ z)$ and $(x_{1},\ y_{1},\ z_{1})$, we see that
\[
(\xi,\ \eta,\ \zeta\,\mid\, x,\ y,\ z) = (\xi_{1},\ \eta_{1},\ \zeta_{1}\,\mid\, x_{1},\ y_{1},\ z_{1}),
\]
and multiplying the first of the relations between two quadrics by an indeterminate quantity $\lambda$, and adding it to the second, we have
\[
(\lambda a + A,\ldots \mid x,\ y,\ z)^{2} = (\lambda a_{1} + A_{1},\ldots \mid x_{1},\ y_{1},\ z_{1})^{2}.
\]

We have thus a linear function and a quadric transformed into functions of the same form by means of the linear substitutions, and any invariant of the system will remain unaltered to a factor pr\`es, such factor being a power of the determinant of substitution. The invariants are, $1^{\circ}$ the discriminant of the quadric; $2^{\circ}$ the reciprocant, considered not as a contravariant of the quadric, but as an invariant of the system. And if we write
\[
K = \mathrm{Disc.}\,(\lambda a + A,\ldots \mid x,\ y,\ z)^{2},
\]
\[
(\mathfrak{A},\ \mathfrak{B},\ \mathfrak{C},\ \mathfrak{F},\ \mathfrak{G},\ \mathfrak{H}\,\mid\,\xi,\ \eta,\ \zeta)^{2} = \mathrm{Recip.}\,(\lambda a + A,\ldots \mid x,\ y,\ z)^{2},
\]
then $K_{1}$, \&c.\ being the analogous expressions for the transformed functions, and the determinant of substitution being represented by $\Pi$, we have
\[
K_{1} = \Pi^{2} K,
\]
\[
(\mathfrak{A}_{1},\ldots \mid \xi_{1},\ \eta_{1},\ \zeta_{1})^{2} = \Pi^{2} (\mathfrak{A},\ldots \mid \xi,\ \eta,\ \zeta)^{2},
\]
and substituting for $\xi_{1}$, $\eta_{1}$, $\zeta_{1}$ their values in terms of $\xi$, $\eta$, $\zeta$, the last equation breaks up into six equations, and we have
\[
K_{1} = \Pi^{2} K,
\]
\[
(\mathfrak{A}_{1},\ldots \mid \alpha,\ \alpha',\ \alpha'')^{2} = \Pi^{2} \mathfrak{A},
\]
\[
(\mathfrak{A}_{1},\ldots \mid \beta,\ \beta',\ \beta'')(\gamma,\ \gamma',\ \gamma'') = \Pi^{2} \mathfrak{F},
\]
which is the system obtained in a somewhat different manner in my former paper.

Putting $f_{1} = g_{1} = h_{1} = F_{1} = G_{1} = H_{1} = 0$, and writing also (which is no additional loss of generality) $a_{1} = b_{1} = c_{1} = 1$, the formulae become
\[
(a,\ b,\ c,\ f,\ g,\ h \mid x,\ y,\ z)^{2} = (1,\ 1,\ 1\,\mid\, x_{1}^{2},\ y_{1}^{2},\ z_{1}^{2}),
\]
\[
(A,\ B,\ C,\ F,\ G,\ H \mid x,\ y,\ z)^{2} = (A_{1},\ B_{1},\ C_{1}\mid x_{1}^{2},\ y_{1}^{2},\ z_{1}^{2}),
\]
viz.\ there are two given quadrics which are to be by the same linear substitution transformed, one of them into the form $x_{1}^{2} + y_{1}^{2} + z_{1}^{2}$ and the other into the form $A_{1} x_{1}^{2} + B_{1} y_{1}^{2} + C_{1} z_{1}^{2}$, where $A_{1}$, $B_{1}$, $C_{1}$ have to be determined. The solution is contained in the following system of formulae, viz.
\[
(A_{1} + \lambda)(B_{1} + \lambda)(C_{1} + \lambda) = \Pi^{2}\cdot\mathrm{Disc.}\,(\lambda a + A,\ldots),
\]

\newpage

%===============================================================
% PDF p.153 = book p.131
%===============================================================
\begin{center}
\small 188]\hfill\textsc{homogeneous functions of the second order.}\hfill\textsc{131}
\end{center}

\medskip

\noindent which gives $A_{1}$, $B_{1}$, $C_{1}$ as the roots of a cubic equation, and gives also
\[
1 = \Pi^{2}\cdot\mathrm{Disc.}\,(a,\ldots) = \Pi^{2}\,\kappa, \text{ or } \Pi^{2} = \frac{1}{\kappa} \text{ suppose,}
\]
and we have then, writing for shortness, $(*\!\!\mid\! X,\ Y,\ Z)$ for
\[
((B_{1} + \lambda)(C_{1} + \lambda),\ (C_{1} + \lambda)(A_{1} + \lambda),\ (A_{1} + \lambda)(B_{1} + \lambda)\mid X,\ Y,\ Z),
\]
\[
(*\!\!\mid\!\alpha^{2},\ \alpha'^{2},\ \alpha''^{2}) = \frac{1}{\kappa}\mathfrak{A},
\]
\[
(*\!\!\mid\!\beta^{2},\ \beta'^{2},\ \beta''^{2}) = \frac{1}{\kappa}\mathfrak{B},
\]
\[
(*\!\!\mid\!\gamma^{2},\ \gamma'^{2},\ \gamma''^{2}) = \frac{1}{\kappa}\mathfrak{C},
\]
\[
(*\!\!\mid\!\beta\gamma,\ \beta'\gamma',\ \beta''\gamma'') = \frac{1}{\kappa}\mathfrak{F},
\]
\[
(*\!\!\mid\!\gamma\alpha,\ \gamma'\alpha',\ \gamma''\alpha'') = \frac{1}{\kappa}\mathfrak{G},
\]
\[
(*\!\!\mid\!\alpha\beta,\ \alpha'\beta',\ \alpha''\beta'') = \frac{1}{\kappa}\mathfrak{H},
\]
where $(\mathfrak{A},\ \mathfrak{B},\ \mathfrak{C},\ \mathfrak{F},\ \mathfrak{G},\ \mathfrak{H})$ are the coefficients of the reciprocant of $(\lambda a + A,\ldots\mid x,\ y,\ z)^{2}$. Writing $\lambda = -A_{1}$, $-B_{1}$, or $-C_{1}$, the quadric functions on the left-hand side become mere monomials, and we have the actual values of the squares and products $\alpha^{2}$, $\beta^{2}$, $\&\text{c.}$ of the coefficients of the linear substitutions: thus $\alpha^{2}$, $\beta^{2}$, $\gamma^{2}$, $\beta\gamma$, $\gamma\alpha$, $\alpha\beta$ are respectively equal to $\mathfrak{A}$, $\mathfrak{B}$, $\mathfrak{C}$, $\mathfrak{F}$, $\mathfrak{G}$, $\mathfrak{H}$, each into the common factor
\[
\frac{1}{\kappa}\,(B_{1} - A_{1})(C_{1} - A_{1}),
\]
the suffix denoting that we are to write in the expressions for $\mathfrak{A}$, $\mathfrak{B}$, $\mathfrak{C}$, $\mathfrak{F}$, $\mathfrak{G}$, $\mathfrak{H}$ the value $-A_{1}$ for $\lambda$; and similarly for the sets $(\alpha',\ \beta',\ \gamma')$ and $(\alpha'',\ \beta'',\ \gamma'')$.

\medskip

\noindent \emph{2, Stone Buildings, 27th March}, 1857.

\newpage

%===============================================================
% PDF p.154 = book p.132 -- Paper 189 begins
%===============================================================
\begin{center}
\small 132\hfill\textsc{}\hfill[189
\end{center}

\bigskip

\begin{center}
{\Large\bfseries 189.}\\[8pt]
{\large NOTE ON A FORMULA IN FINITE DIFFERENCES.}\\[6pt]
\small [From the \emph{Quarterly Mathematical Journal}, vol.~\textsc{ii}.\ (1858), pp.~198--201.]
\end{center}

\bigskip

\noindent \textsc{In} Jacobi's Memoir ``De usu legitimo formulae summatoriae Maclaurinianae,'' \emph{Crelle}, t.~\textsc{xii}.\ [1834], pp.~263--273 (1834), expressions are given for the sums of the odd powers of the natural numbers $1$, $2$, $3 \ldots x$ in terms of the quantity
\[
u = x(x+1),
\]
viz.\ putting for shortness
\[
Sx^{r} = 1^{r} + 2^{r} + \ldots + x^{r},
\]
the expressions in question are
\begin{align*}
Sx^{3} &= \tfrac{1}{4} u^{2}, \\
Sx^{5} &= \tfrac{1}{6} u^{3}\,(u - \tfrac{1}{3}), \\
Sx^{7} &= \tfrac{1}{8} u^{4}\,(u^{2} - \tfrac{4}{3} u + \tfrac{2}{3}), \\
Sx^{9} &= \tfrac{1}{10} u^{2}\,(u^{4} - 3 u^{3} + \tfrac{17}{5} u^{2} - \tfrac{10}{3} u + \tfrac{4}{3}), \\
Sx^{11} &= \tfrac{1}{12} u^{2}\,(u^{5} - \tfrac{14}{3} u^{4} + \tfrac{32}{3} u^{3} - \tfrac{72}{5} u^{2} + \tfrac{208}{15} u - \tfrac{32}{5}), \\
\&\text{c.,}
\end{align*}
which, especially as regards the lower powers, are more simple than the ordinary expressions in terms of $x$.

The expressions are continued by means of a recurring formula, viz.\ if
\[
Sx^{2p-1} = \frac{1}{2p - 2}\left[a_{0} u^{p+1} - a_{1} u^{p+2} + \ldots + (-)^{p-1} a_{p-1} u^{2}\right],
\]
\[
Sx^{2p+1} = \frac{1}{2p}\left[u^{p} - b_{1} u^{p-1} + \ldots + (-)^{p}\, b_{p}\, u^{0}\right],
\]

\newpage

%===============================================================
% PDF p.155 = book p.133
%===============================================================
\begin{center}
\small 189]\hfill\textsc{note on a formula in finite differences.}\hfill\textsc{133}
\end{center}

\medskip

\noindent then
\begin{align*}
2p(2p-1)\,a_{1} &= (2p-2)(2p-3)\,b_{1} - p(p-1), \\
2p(2p-1)\,a_{2} &= (2p-4)(2p-5)\,b_{2} - (p-1)(p-2)\,b_{1}, \\
2p(2p-1)\,a_{3} &= (2p-6)(2p-7)\,b_{3} - (p-2)(p-3)\,b_{2}, \\
&\hspace{2em} \vdots \\
2p(2p-1)\,a_{p-2} &= 5\cdot 6\,b_{p-2} - 3\cdot 4\,b_{p-3}, \\
0 &= 3\cdot 4\,b_{p-1} - 2\cdot 3\,b_{p-2},
\end{align*}
by means of which the coefficients $b$ can be determined when the coefficients $a$ are known.

Jacobi remarks also that the expressions for the sums of the even powers may be obtained from those for the odd powers by means of the formula
\[
Sx^{2p} = \frac{1}{2p+2}\,\tfrac{d}{du}\,Sx^{2p+1},
\]
which shows that any such sum will be of the form $(2x+1)u$ into a rational and integral function of $u$: thus in particular
\[
Sx^{2} = \tfrac{1}{3}(2x+1)\,u.
\]

To show \emph{a priori} that $Sx^{2p+1}$ can be expressed as a rational and integral function of $u$, it may be remarked that $Sx^{2p+1} = \phi_{1}x$ where $\phi_{1}x$ denotes the summatory integral $\Sigma(x+1)^{2p+1}$, taken so as to vanish for $x = 0$: $\phi_{1}x$ is a rational and integral function of $x$ of the degree $2p+2$, and which, as is well known, contains $x^{2}$ as a factor. Suppose that $y$ is any positive or negative integer less than $x$, we have
\[
\phi_{1}x - \phi_{1}y = (y+1)^{2p+1} + (y+2)^{2p+1} + \ldots + x^{2p+1},
\]
and in particular putting $y = -1 - x$,
\[
\phi_{1}x - \phi_{1}(-1-x) = (-x)^{2p+1} + (1-x)^{2p+1} + \ldots + x^{2p+1} = 0,
\]
since the terms destroy each other in pairs; we have therefore $\phi_{1}x = \phi_{1}(-1-x)$. Now $u = x^{2} + x$, or writing this equation under the form $x^{2} = -x + u$, we see that any rational and integral function of $x$ may be reduced to the form $Px + Q$, where $P$ and $Q$ are rational and integral functions of $u$. Write therefore $\phi_{1}x = Px + Q$: the substitution of $-1-x$ in the place of $x$ leaves $u$ unaltered, and the equation $\phi_{1}x = \phi_{1}(-1-x)$ thus shows that $P = 0$; we have therefore $\phi_{1}x = Q$, a rational and integral function of $u$. Moreover $\phi_{1}x$ as containing the factor $x^{2}$, must clearly contain the factor $u^{2}$, and the expressions for $Sx^{2p+1}$ are thus shown to be of the form given by Jacobi.

We may obtain a finite expression for $Sx^{n}$ in terms of the differences of $0^{n}$ as follows: we have
\[
Sx^{n} = 1^{n} + 2^{n} + \ldots + x^{n} = \left\{(1+\Delta) + (1+\Delta)^{2} + \ldots + (1+\Delta)^{x}\right\}\!0^{n} = \frac{1+\Delta}{\Delta}\,\{(1+\Delta)^{x} - 1\}\,0^{n},
\]

\newpage

%===============================================================
% PDF p.156 = book p.134
%===============================================================
\begin{center}
\small 134\hfill\textsc{note on a formula in finite differences.}\hfill[189
\end{center}

\medskip

\noindent and putting $(1+\Delta)^{x} = e^{x\log(1+\Delta)}$ and observing that the term independent of $x$ vanishes, and that the terms containing powers higher than $x^{n+1}$ also vanish, we have
\[
Sx^{n} = S_{k}\left\{\frac{1+\Delta}{\Delta}\log^{k}(1+\Delta)\right\}0^{n}\cdot\frac{x^{k}}{\Pi k},
\]
where the summation with respect to $k$, extends from $k = 1$ to $k = n + 1$, or what is the same thing (since the term corresponding to $k = 1$ in fact vanishes) from $k = 2$ to $k = n + 1$.

The equation $x^{2} = -x + u$ gives
\[
x^{k} = P_{k} x + Q_{k},
\]
and it is easy to see that writing for shortness
\[
M_{k} = 1 + \frac{k-3}{1}\,u + \frac{k-4\cdot k-5}{1\cdot 2}\,u^{2} + \frac{k-5\cdot k-6\cdot k-7}{1\cdot 2\cdot 3}\,u^{3} + \ldots,
\]
where the series is to be continued to the term $u^{\frac{1}{2}(k-2)}$ or $u^{\frac{1}{2}(k-3)}$ according as $k$ is even or odd, we have
\[
P_{k} = (-)^{k+1}\,M_{k+1}, \quad Q_{k} = (-)^{k}\,u\,M_{k},
\]
we have consequently
\[
Sx^{n} = x S_{k}\left\{\frac{1+\Delta}{\Delta}\log^{k}(1+\Delta)\right\}0^{n}\cdot\frac{(-)^{k+1}\,M_{k+1}}{\Pi k}
+ S_{k}\left\{\frac{1+\Delta}{\Delta}\log^{k}(1+\Delta)\right\}0^{n}\cdot\frac{(-)^{k}\,u\,M_{k}}{\Pi k}.
\]

If $n$ is odd, $= 2p + 1$, then (by what precedes) the first term vanishes, or we have
\[
S_{k}\left\{\frac{1+\Delta}{\Delta}\log^{k}(1+\Delta)\right\}0^{2p+1}\cdot\frac{(-)^{k+1}\,M_{k+1}}{\Pi k} = 0,\quad (k = 1 \text{ to } k = 2p + 2),
\]
and the formula becomes
\[
Sx^{2p+1} = S_{k}\left\{\frac{1+\Delta}{\Delta}\log^{k}(1+\Delta)\right\}0^{2p+1}\cdot\frac{(-)^{k}\,u\,M_{k}}{\Pi k},\quad (k = 1 \text{ to } k = 2p + 2),
\]
which it may be noticed puts in evidence the factor $u$ but not the factor $u^{2}$.

If $n$ is even, $= 2p$, then (by what precedes) the coefficient of $x$ is to the constant term in the ratio $2 : 1$, or we have
\[
S_{k}\left\{\frac{1+\Delta}{\Delta}\log^{k}(1+\Delta)\right\}0^{2p}\cdot\frac{(-)^{k+1}\,(M_{k+1} - 2 u M_{k})}{\Pi k} = 0,\ (k = 1 \text{ to } k = 2p + 1),
\]
and the formula becomes
\[
Sx^{2p} = (2x+1)\,S_{k}\left\{\frac{1+\Delta}{\Delta}\log^{k}(1+\Delta)\right\}0^{2p}\cdot\frac{(-)^{k}\,u\,M_{k}}{\Pi k},\ (k = 1 \text{ to } k = 2p + 1).
\]

\newpage

%===============================================================
% PDF p.157 = book p.135
%===============================================================
\begin{center}
\small 189]\hfill\textsc{note on a formula in finite differences.}\hfill\textsc{135}
\end{center}

\medskip

\noindent The values of the functions $M$ are as follows:
\begin{align*}
M_{1} &= 0, \\
M_{2} &= 1, \\
M_{3} &= 1, \\
M_{4} &= 1 + u, \\
M_{5} &= 1 + 2u, \\
M_{6} &= 1 + 3u + u^{2}, \\
M_{7} &= 1 + 4u + 3u^{2}, \\
&\&\text{c.}
\end{align*}

As a simple example of the formulae, we have
\begin{align*}
Sx^{3} &= \left\{\frac{1+\Delta}{\Delta}\log^{2}(1+\Delta)\right\}0^{3}\cdot\tfrac{1}{2}\,u \\
&\hspace{2em} + \left\{\frac{1+\Delta}{\Delta}\log^{3}(1+\Delta)\right\}0^{3}\cdot -\tfrac{1}{6}\,u \\
&\hspace{2em} + \left\{\frac{1+\Delta}{\Delta}\log^{4}(1+\Delta)\right\}0^{3}\cdot\tfrac{1}{24}\,(u + u^{2}),
\end{align*}
and the coefficients are
\begin{align*}
(\Delta - \tfrac{1}{12}\Delta^{3})\,0^{3} &= 1 - \tfrac{1}{4}\cdot 6 = \tfrac{1}{2}, \\
(\Delta^{2} - \tfrac{1}{2}\Delta^{3})\,0^{3} &= 6 - \tfrac{1}{2}\cdot 6 = 3, \\
\Delta^{3}\,0^{3} &= 6,
\end{align*}
and therefore
\[
Sx^{3} = \tfrac{1}{2} u - \tfrac{1}{2} u + \tfrac{1}{4}\,(u + u^{2}) = \tfrac{1}{4} u^{2},
\]
which is right; the example shows however that the calculation for the higher powers would be effected more readily by means of Jacobi's recurring formula.

\medskip

\noindent \emph{2, Stone Buildings, 27th Oct.}, 1857.

\newpage

%===============================================================
% PDF p.158 = book p.136 -- Paper 190 begins
%===============================================================
\begin{center}
\small 136\hfill\textsc{}\hfill[190
\end{center}

\bigskip

\begin{center}
{\Large\bfseries 190.}\\[8pt]
{\large ON THE SYSTEM OF CONICS WHICH PASS THROUGH THE\\
SAME FOUR POINTS.}\\[6pt]
\small [From the \emph{Quarterly Mathematical Journal}, vol.~\textsc{ii}.\ (1858), pp.~206--207.]
\end{center}

\bigskip

\noindent \textsc{I consider} the system of conics passing through the same four points; these points may be real or imaginary, but it is assumed that there is a real system of conics, this will in fact be the case if two conics of the system are real. The four points are therefore given as the points of intersection of two real conics, and it will be proper to assume in the first instance that the conics intersect in four separate and distinct points, none of them at infinity. The four points may be all real, or two real and two imaginary, or all imaginary.

First, if the points are all real, we have here two cases, viz.\ each of the points may lie outside of the triangle formed by the other three, or as this may be expressed, the points may form a convex quadrangle; or else one of the points may be inside the triangle formed by the other three, or as this may be expressed, the points may form a triangle and interior point. In each case the pairs of lines joining the points, two and two together, will be conics (degenerate hyperbolas) forming part of the system of conics. Consider the two cases separately.

\medskip

\begin{center}
[\emph{Fig.\ A.\ omitted: four real points forming a convex quadrangle; six lines joined and labelled.}]
\end{center}

\newpage

%===============================================================
% PDF p.159 = book p.137
%===============================================================
\begin{center}
\small 190]\hfill\textsc{on the system of conics \&c.}\hfill\textsc{137}
\end{center}

\medskip

\noindent \emph{Fig.~A}. Four real points forming a convex quadrangle. The system contains two parabolas, and the pairs of lines and the parabolas divide the plane of the figure into five distinct regions, one of which contains only ellipses, and the other four contain each of them hyperbolas.

\emph{Fig.~A$'$}. Four real points forming a triangle and interior point. The system does not contain any parabolas, the three pairs of lines divide the plane of the figure into three distinct regions, each of which contains only hyperbolas.

\begin{center}
[\emph{Fig.\ A$'$.\ omitted.}]
\end{center}

Next, if the points are two of them real and two of them imaginary. The line joining the two imaginary points will be real and this line may meet the line joining the two real points, in a point outside the two real points, or included between them, i.e.\ the real centre of the quadrangle may lie outside the real points, or may be included between them; I consider the two cases separately.

\emph{Fig.~B}. Two real and two imaginary points, the real centre of the quadrangle lying outside the real points. The system contains two parabolas, and these with the line joining the two real points and the line joining the two imaginary points divide the plane of the figure into three regions, one of which contains ellipses and the other two contain each of them hyperbolas.

\begin{center}
[\emph{Fig.\ B.\ omitted.}]
\end{center}

\newpage

%===============================================================
% PDF p.160 = book p.138
%===============================================================
\begin{center}
\small 138\hfill\textsc{on the system of conics which pass through the same four points.}\hfill[190
\end{center}

\medskip

\noindent \emph{Fig.~B$'$}. Two real and two imaginary points, the real centre of the quadrangle lying between the real points. There are no parabolas, and the system contains only hyperbolas.

\begin{center}
[\emph{Fig.\ B$'$.\ omitted.}]
\end{center}

Lastly, when the four points are imaginary. We have here only a single case.

\emph{Fig.~C}. Four imaginary points. The points lie on two real lines, there are (besides the point of intersection of these lines) two other real centres of the quadrangle, which lie harmonically with respect to the two lines. The system contains two parabolas and these and the two lines divide the plane of the figure into four regions, two of which contain each of them ellipses, and the other two contain each of them hyperbolas.

\begin{center}
[\emph{Fig.\ C.\ omitted.}]
\end{center}

\newpage

%===============================================================
% PDF p.161 = book p.139 -- Paper 191 begins
%===============================================================
\begin{center}
\small 191]\hfill\textsc{}\hfill\textsc{139}
\end{center}

\bigskip

\begin{center}
{\Large\bfseries 191.}\\[8pt]
{\large NOTE ON THE EXPANSION OF THE TRUE ANOMALY.}\\[6pt]
\small [From the \emph{Quarterly Mathematical Journal}, vol.~\textsc{ii}.\ (1858), pp.~229--232.]
\end{center}

\bigskip

\noindent \textsc{If} the true anomaly and the mean anomaly are respectively denoted by $u$, $m$, and if $e$ be the eccentricity, then as usual $u - e\sin u = m$; and if we write
\[
\lambda = \frac{1 - \sqrt{(1 - e^{2})}}{e},
\]
and take $c$ to denote the base of the hyperbolic system of logarithms, we have
\[
u = m + 2\Sigma_{1}^{\infty} A_{r}\,\frac{\sin r m}{r},
\]
and
\[
A_{r} = \lambda^{r}\,c^{-\frac{1}{2}r(\lambda - \lambda^{-1})}\, + \lambda^{-r}\,c^{\frac{1}{2}r(\lambda - \lambda^{-1})},
\]
where, after expanding the exponentials, the negative powers of $\lambda$ are to be rejected and the term independent of $\lambda$ is to be multiplied by $\tfrac{1}{2}$ (see \emph{Camb.\ Math.\ Journal}, t.~\textsc{i}.\ [1839] p.~228 and t.~\textsc{iii}.\ [1843] p.~165, [4]).

It is easily seen that $e^{r}$ is the lowest power of $e$ which enters into the value of $A_{r}$, and the question arises to find the numerical coefficient of the term in question; this is readily obtained from the formula; in fact considering first a term of the form
\[
\lambda^{-r}\,c^{s}\,(\lambda - \lambda^{-1})^{s},
\]
since $\lambda$ is itself of the order $e$, when the negative powers of $\lambda$ are rejected this is at least of the order $e^{s}$ and it is consequently to be neglected if $s > r$. But if $s < r$ all the powers of $\lambda$ are negative and the term is to be rejected. The only case to be

\newpage

%===============================================================
% PDF p.162 = book p.140
%===============================================================
\begin{center}
\small 140\hfill\textsc{note on the expansion of the true anomaly.}\hfill[191
\end{center}

\medskip

\noindent considered is therefore that of $s = r$, in which case there is a term containing $e^{r}$. We thus obtain from $\lambda^{-r}\,c^{\frac{1}{2}r(\lambda - \lambda^{-1})}$ the term
\[
\frac{1}{2^{r}}\cdot\frac{r^{r} e^{r}}{1\cdot 2\cdot 3\ldots r}.
\]

In the next place a term of the form $\lambda^{r}\,c^{s}\,(\lambda - \lambda^{-1})^{s}$ is at least of the order $e^{r}$ if $s > r$, or the terms to be considered are those for which $s = r$ or $r < r$. But in such term the only part of the order $e^{r}$ is
\[
(-)^{s}\,2^{-r+s}\,e^{r},
\]
or, since neglecting higher powers of $e$ we have $\lambda = \tfrac{1}{2} e$, this is
\[
(-)^{s}\,2^{-r+s}\,e^{r},
\]
and the set of terms arising from
\[
\lambda^{r}\,c^{-\frac{1}{2}r(\lambda - \lambda^{-1})}
\]
is
\[
\frac{e^{r}}{2^{r}}\left\{1 + \frac{r}{1} + \frac{r^{2}}{1\cdot 2} + \ldots + \frac{r^{r-1}}{1\cdot 2\ldots(r - 1)} + \tfrac{1}{2}\,\frac{r^{r}}{1\cdot 2\ldots r}\right\},
\]
the last term being divided by $2$ because arising from a term independent of $\lambda$. Hence the first term of $A_{r}$ is
\[
\frac{e^{r}}{2^{r}}\left\{1 + \frac{r}{1} + \frac{r^{2}}{1\cdot 2} + \ldots + \frac{r^{r}}{1\cdot 2\ldots r}\right\},
\]
a result which it may be remarked is contained in the general formula given in Hansen's Memoir ``Entwickelung des Products u.s.w.,'' \emph{Leipzig Trans.}, t.~\textsc{ii}.\ p.~277 (1853).

The preceding expression is
\[
\frac{e^{r}}{2^{r}}\,c^{r}\,\frac{1}{\Gamma(r+1)}\int_{r}^{\infty}\,x^{r} c^{-x}\,dx,
\]
and to find its value when $r$ is large, we have
\begin{align*}
\int_{r}^{\infty}\,x^{r} c^{-x}\,dx &= \int_{0}^{\infty}\,(y + r)^{r}\,c^{-y - r}\,dy = r^{r} c^{-r}\int_{0}^{\infty}\,\left(1 + \tfrac{y}{r}\right)^{r}\,c^{-y}\,dy \\
&= r^{r} c^{-r}\int_{0}^{\infty}\,c^{-y + r\log(1 + \frac{y}{r})}\,dy \\
&= r^{r} c^{-r}\int_{0}^{\infty}\,c^{-\frac{y^{2}}{2r} + \frac{y^{3}}{3r^{2}} - \&\text{c.}}\,dy \\
&= r^{r} c^{-r}\int_{0}^{\infty}\,\left(1 + \frac{y^{3}}{3r^{2}} + \ldots\right)\,c^{-\frac{y^{2}}{2r}}\,dy \\
&= r^{r} c^{-r}\,\sqrt{2 r}\int_{0}^{\infty}\,\left(1 + \frac{2\sqrt{2}}{3\sqrt{r}}\,z^{3} + \ldots\right)\,c^{-z^{2}}\,dz,
\end{align*}

\newpage

%===============================================================
% PDF p.163 = book p.141
%===============================================================
\begin{center}
\small 191]\hfill\textsc{note on the expansion of the true anomaly.}\hfill\textsc{141}
\end{center}

\medskip

\noindent or neglecting all the terms except the first, this is
\[
= r^{r} c^{-r}\,\sqrt{2 r}\,\int_{0}^{\infty}\,c^{-z^{2}}\,dz = \sqrt{2 \pi r}\,r^{r}\,c^{-r}.
\]

Hence multiplying by $\dfrac{1}{2^{r}}\,e^{r}\,c^{r}\,\dfrac{1}{\Gamma(r+1)}$ and observing that when $r$ is large, we have, by a well-known formula,
\[
\Gamma(r+1) = \sqrt{2 \pi r}\,r^{r}\,c^{-r},
\]
we obtain finally the result that when $r$ is large the first term of $A_{r}$ is approximately
\[
= \left(\frac{e c}{2}\right)^{r}.
\]

I take the opportunity of mentioning the following somewhat singular theorem, which seems to belong to a more general theory: viz.\ if $u - e\sin u = m$, then we have
\[
\log(1 - e\cos u) = \frac{1}{a}\log(1 - a e\cos\phi),
\]
where
\[
\phi - \frac{1}{a}\tan\phi = m,
\]
provided that the negative powers of $a$ are rejected, and $a$ is then put equal to unity.

To show this, we have by Lagrange's theorem, observing that
\[
\frac{d}{dm} F(1 - e\cos m) = e\sin m\,F'(1 - e\cos m),
\]
\begin{align*}
F(1 - e\cos u) &= F(1 - e\cos m) + \frac{e^{2}}{1\cdot 2}\,\frac{d}{dm}\sin^{2} m\,F'(1 - e\cos m) \\
&\hspace{2em} + \frac{e^{3}}{1\cdot 2\cdot 3}\,\frac{d^{2}}{dm^{2}}\sin^{3} m\,F''(1 - e\cos m) + \&\text{c.,}
\end{align*}
and the coefficient of $e^{r}$ in $F(1 - e\cos u)$ is
\[
\frac{(-)^{r}}{1\cdot 2\ldots(r - 1)}\left\{\tfrac{1}{r}\,F_{r}\cos^{r} m + \frac{r-1}{1}\,F_{r-1}\cos^{r-1} m\sin^{2} m - \frac{(r-1)(r-2)}{1\cdot 2}\,F_{r-2}\,\tfrac{d}{dm}(\cos^{r-3} m\sin^{3} m) + \&\text{c.}\right\},
\]
where $F_{r} = F^{(r)}(1)$.

\newpage

%===============================================================
% PDF p.164 = book p.142
%===============================================================
\begin{center}
\small 142\hfill\textsc{note on the expansion of the true anomaly.}\hfill[191
\end{center}

\medskip

\noindent Hence in particular when $F x = \log x$, $F_{r} = (-)^{r-1}\cdot 1\cdot 2\ldots(r-1)$ and thence the coefficient of $e^{r}$ in $\log(1 - e\cos u)$ is
\[
- \left\{\tfrac{1}{r}\cos^{r} m - \frac{1}{1\cdot 2}\cos^{r-2} m\sin^{2} m - \frac{1}{1\cdot 2}\,\frac{d}{dm}(\cos^{r-3} m\sin^{3} m) - \&\text{c.}\right\},
\]
continued as long as the exponent of $\cos m$ is not negative. Now in the expansion of $\frac{1}{a}\log(1 - a e\cos\phi)$, where $\phi - \frac{1}{a}\tan\phi = m$, the coefficient of $e^{r}$ is $-\frac{1}{r}\,a^{r-1}\cos^{r}\phi$, and this (by Lagrange's theorem) is equal to
\[
- \frac{1}{r}\,a^{r-1}\left\{\cos^{r} m - \frac{1}{1\cdot 2}\,r\cos^{r-1} m\sin m\tan m - \frac{1}{1\cdot 2\cdot a}\,\frac{d}{dm}(r\cos^{r-2} m\sin^{2} m\tan^{2} m) - \&\text{c.}\right\}
\]
\[
= -\left\{\tfrac{1}{r}\,a^{r-1}\cos^{r} m - \tfrac{1}{1}\,a^{r-2}\cos^{r-2} m\sin^{2} m - \frac{1}{1\cdot 2}\,a^{r-3}\cos^{r-3} m\sin^{2} m - \&\text{c.}\right\},
\]
where the series is continued indefinitely; but if we reject the negative powers of $a$ and then put $a$ equal to unity this is precisely equal to the former expression for the coefficient of $e^{r}$, and the formula is thus shown to be true.

\medskip

\noindent \emph{2, Stone Buildings, W.C., 17th Nov.}, 1857.

\newpage

%===============================================================
% PDF p.165 = book p.143 -- Paper 192 begins
%===============================================================
\begin{center}
\small 192]\hfill\textsc{}\hfill\textsc{143}
\end{center}

\bigskip

\begin{center}
{\Large\bfseries 192.}\\[8pt]
{\large ON THE AREA OF THE CONIC SECTION REPRESENTED BY\\
THE GENERAL TRILINEAR EQUATION OF THE SECOND\\
DEGREE.}\\[6pt]
\small [From the \emph{Quarterly Mathematical Journal}, vol.~\textsc{ii}.\ (1858), pp.~248--253.]
\end{center}

\bigskip

\noindent [\textsc{The} original title was ``Direct Investigation of the Question discussed in the Foregoing Paper,'' viz.\ a Paper with the present title by N.~M.~Ferrers (now Dr Ferrers), pp.~247--248. The area $S$ of the conic section represented by the general equation $(A,\ B,\ C,\ A',\ B',\ C'\,\mid\, x,\ y,\ z)^{2} = 1$, where the coordinates are connected by the equation $x + y + z = 1$, was by considerations founded on the form of the function found to be
\[
S = \frac{2\pi\,(AA'^{2} + BB'^{2} + CC'^{2} - ABC - 2A'B'C')\,\Delta}{[A'^{2} - BC + B'^{2} - CA + C'^{2} - AB + 2(B'C' - AA') + 2(C'A' - BB') + 2(A'B' - CC')]^{\frac{3}{2}}},
\]
where $\Delta$ is the area of the fundamental triangle; and it was remarked that a similar method might be applied to determine the area of the conic section when it is defined by the distances of its several tangents from three given points.]

The position of a point $P$ being determined as in the foregoing paper, let $\alpha$, $\beta$, $\gamma$ denote in like manner the coordinates of a point $O$, we have
\[
\alpha + \beta + \gamma = 1,
\]
and consequently if $\xi$, $\eta$, $\zeta$ are the relative coordinates $x - \alpha$, $y - \beta$, $z - \gamma$, we have
\[
\xi + \eta + \zeta = 0.
\]

\newpage

%===============================================================
% PDF p.166 = book p.144
%===============================================================
\begin{center}
\small 144\hfill\textsc{on the area of the conic section represented by the}\hfill[192
\end{center}

\medskip

\noindent The expression for the distance of the two points $O$, $P$ is readily obtained in terms of the relative coordinates, viz.\ calling this distance $r$, we have
\[
r^{2} = L\xi^{2} + M\eta^{2} + N\zeta^{2},
\]
where, if $l$, $m$, $n$ are the sides of the triangle $ABC$, we have
\begin{align*}
L &= \tfrac{1}{2}(m^{2} + n^{2} - l^{2}), \\
M &= \tfrac{1}{2}(n^{2} + l^{2} - m^{2}), \\
N &= \tfrac{1}{2}(l^{2} + m^{2} - n^{2});
\end{align*}
and it is to be remarked that these values give
\[
MN + NL + LM = \tfrac{1}{4}\,(2 m^{2} n^{2} + 2 n^{2} l^{2} + 2 l^{2} m^{2} - l^{4} - m^{4} - n^{4}) = 4\Delta^{2},
\]
if $\Delta$ denote the area of the triangle $ABC$.

Consider now a conic
\[
(a,\ b,\ c,\ f,\ g,\ h\,\mid\, x,\ y,\ z)^{2},
\]
and suppose as usual that $\mathfrak{A}$, $\mathfrak{B}$, $\mathfrak{C}$, $\mathfrak{F}$, $\mathfrak{G}$, $\mathfrak{H}$ are the inverse coefficients and that $K$ is the discriminant, suppose also for shortness
\[
P = (\mathfrak{A},\ \mathfrak{B},\ \mathfrak{C},\ \mathfrak{F},\ \mathfrak{G},\ \mathfrak{H}\,\mid\, 1,\ 1,\ 1)^{2}.
\]
The coordinates of the centre being $\alpha$, $\beta$, $\gamma$, we have
\begin{align*}
\alpha &= \frac{1}{P}\,(\mathfrak{A},\ \mathfrak{H},\ \mathfrak{G}\,\mid\, 1,\ 1,\ 1), \\
\beta &= \frac{1}{P}\,(\mathfrak{H},\ \mathfrak{B},\ \mathfrak{F}\,\mid\, 1,\ 1,\ 1), \\
\gamma &= \frac{1}{P}\,(\mathfrak{G},\ \mathfrak{F},\ \mathfrak{C}\,\mid\, 1,\ 1,\ 1),
\end{align*}
and writing as before $\xi$, $\eta$, $\zeta$ for $x - \alpha$, $y - \beta$, $z - \gamma$, so that $\xi$, $\eta$, $\zeta$ are the coordinates of a point $P$ of the conic, in relation to the centre, we have $x$, $y$, $z$ respectively equal to $\xi + \alpha$, $\eta + \beta$, $\zeta + \gamma$, and the equation of the conic gives
\[
(a,\ldots \mid \xi + \alpha,\ \eta + \beta,\ \zeta + \gamma)^{2} = 0,
\]
which may be written
\begin{align*}
&(a,\ldots \mid \xi,\ \eta,\ \zeta)^{2} \\
&\quad + 2\,(a,\ldots \mid \xi,\ \eta,\ \zeta)\,(\xi,\ \eta,\ \zeta\,\mid\,\alpha,\ \beta,\ \gamma)_{1} \\
&\quad + (a,\ldots \mid \alpha,\ \beta,\ \gamma)^{2} = 0.
\end{align*}

\newpage

%===============================================================
% PDF p.167 = book p.145
%===============================================================
\begin{center}
\small 192]\hfill\textsc{general trilinear equation of the second degree.}\hfill\textsc{145}
\end{center}

\medskip

\noindent Now observing the equations
\begin{align*}
(a,\ h,\ g\,\mid\,\alpha,\ \beta,\ \gamma) &= \tfrac{K}{P}, \\
(h,\ b,\ f\,\mid\,\alpha,\ \beta,\ \gamma) &= \tfrac{K}{P}, \\
(g,\ f,\ c\,\mid\,\alpha,\ \beta,\ \gamma) &= \tfrac{K}{P},
\end{align*}
we have
\[
(a,\ldots \mid \alpha,\ \beta,\ \gamma)\,(\xi,\ \eta,\ \zeta) = \tfrac{K}{P}\,(\xi + \eta + \zeta) = 0,
\]
\[
(a,\ldots \mid \alpha,\ \beta,\ \gamma)^{2} = \tfrac{K}{P}\,(\alpha + \beta + \gamma) = \tfrac{K}{P}.
\]
and the equation of the conic gives therefore
\[
(a,\ldots \mid \xi,\ \eta,\ \zeta)^{2} + \tfrac{K}{P} = 0,
\]
and we have as before
\[
\xi + \eta + \zeta = 0.
\]

To find the axes we have only to make
\[
r^{2} = L\xi^{2} + M\eta^{2} + N\zeta^{2},
\]
a maximum or minimum, $\xi$, $\eta$, $\zeta$ varying subject to the preceding two conditions; this gives
\begin{align*}
(a,\ h,\ g\,\mid\,\xi,\ \eta,\ \zeta) + \lambda L\xi + \mu &= 0, \\
(h,\ b,\ f\,\mid\,\xi,\ \eta,\ \zeta) + \lambda M\eta + \mu &= 0, \\
(g,\ f,\ c\,\mid\,\xi,\ \eta,\ \zeta) + \lambda N\zeta + \mu &= 0,
\end{align*}
and multiplying by $\xi$, $\eta$, $\zeta$ adding and reducing, we have
\[
-\tfrac{K}{P} + \lambda r^{2} = 0,
\]
which gives
\[
\lambda = \tfrac{K}{P r^{2}}.
\]
Substituting this value, and joining to the resulting three equations the equation
\[
\xi + \eta + \zeta = 0.
\]

\newpage

%===============================================================
% PDF p.168 = book p.146
%===============================================================
\begin{center}
\small 146\hfill\textsc{on the area of the conic section represented by}\hfill[192
\end{center}

\medskip

\noindent we may eliminate $\xi$, $\eta$, $\zeta$, $\mu$, and the result is
\[
\begin{vmatrix}
a + \tfrac{K L}{P r^{2}}, & h, & g, & 1 \\
h, & b + \tfrac{K M}{P r^{2}}, & f, & 1 \\
g, & f, & c + \tfrac{K N}{P r^{2}}, & 1 \\
1, & 1, & 1, & \cdot
\end{vmatrix} = 0,
\]
which may also be written
\[
(\mathfrak{A}',\ \mathfrak{B}',\ \mathfrak{C}',\ \mathfrak{F}',\ \mathfrak{G}',\ \mathfrak{H}' \mid 1,\ 1,\ 1)^{2} = 0,
\]
where $(\mathfrak{A}',\ldots)$ are what $(\mathfrak{A},\ldots)$ become when $a$, $b$, $c$ are changed into
\[
a + \tfrac{KL}{P r^{2}}, \quad b + \tfrac{KM}{P r^{2}}, \quad c + \tfrac{KN}{P r^{2}};
\]
we in fact have
\[
\mathfrak{A}' = \mathfrak{A} + \tfrac{K}{P r^{2}}\,(bN + cM) + \tfrac{K^{2}}{P^{2} r^{4}}\,MN,
\]
\[
\vdots
\]
\[
\mathfrak{F}' = \mathfrak{F} - \tfrac{K}{P r^{2}}\,L f,
\]
\[
\vdots
\]
and (observing the value of $P$) the result consequently is
\[
P + \tfrac{K}{P r^{2}}\left\{(b + c - 2 f)L + (c + a - 2 g)M + (a + b - 2 h)N\right\} + \tfrac{K^{2}}{P^{2} r^{4}}\,(MN + NL + LM) = 0,
\]
which may also be written
\[
P r^{4} + P K r^{2}\left\{(b + c - 2 f)L + (c + a - 2 g)M + (a + b - 2 h)N\right\} + 4\Delta^{2} K^{2} = 0.
\]
Hence if $r_{1}$, $r_{2}$ are the two semiaxes, we have
\[
r_{1}^{2} r_{2}^{2} = \frac{4\Delta^{2} K^{2}}{P^{2}},
\]
and the area is $\pi r_{1} r_{2}$, which is equal to
\[
\frac{2\pi K\Delta}{\sqrt{(P^{3})}},
\]
which agrees with Mr Ferrers' result.

The formula $r^{2} = L\xi^{2} + M\eta^{2} + N\zeta^{2}$ which is assumed in the preceding investigation may be proved as follows:

\newpage

%===============================================================
% PDF p.169 = book p.147
%===============================================================
\begin{center}
\small 192]\hfill\textsc{general trilinear equation of the second degree.}\hfill\textsc{147}
\end{center}

\medskip

\noindent Writing $a$, $b$, $c$ (instead of $l$, $m$, $n$) for the sides of the fundamental triangle and $A$, $B$, $C$ for the angles, the equation in question is
\[
r^{2} = b c\cos A\cdot\xi^{2} + c a\cos B\,\eta^{2} + a b\cos C\,\zeta^{2}.
\]
Now writing $\alpha$, $\beta$, $\gamma$ for the inclinations of the line $r$ to the sides of the triangle, we have
\[
A = \beta - \gamma,
\]
\[
B = \gamma - \alpha,
\]
\[
C = \pi + \alpha - \beta.
\]
Moreover taking for a moment $\lambda$, $\mu$, $\nu$ to denote the perpendiculars from the angles on the opposite sides, we have
\[
\lambda = c\sin B = b\sin C,
\]
\[
\mu = a\sin C = c\sin A,
\]
\[
\nu = b\sin A = a\sin B,
\]
and
\[
\xi = \frac{r\sin\alpha}{\lambda},\quad \eta = \frac{r\sin\beta}{\mu},\quad \zeta = \frac{r\sin\gamma}{\nu};
\]
the values of $\xi^{2}$, $\eta^{2}$, $\zeta^{2}$ consequently are
\[
\frac{r^{2}\sin^{2}\alpha}{b c\sin B\sin C},\quad \frac{r^{2}\sin^{2}\beta}{c a\sin C\sin A},\quad \frac{r^{2}\sin^{2}\gamma}{a b\sin A\sin B},
\]
and the equation to be proved becomes
\[
1 = \frac{\cos A\sin^{2}\alpha}{\sin B\sin C} + \frac{\cos B\sin^{2}\beta}{\sin C\sin A} + \frac{\cos C\sin^{2}\gamma}{\sin A\sin B},
\]
or, what is the same thing,
\[
\sin A\sin B\sin C = \sin A\cos A\sin^{2}\alpha + \sin B\cos B\sin^{2}\beta + \sin C\cos C\sin^{2}\gamma,
\]
or again
\[
4\sin A\sin B\sin C = \sin 2A\,(1 - \cos 2\alpha) + \sin 2B\,(1 - \cos 2\beta) + 2C\,(1 - \cos 2\gamma),
\]
or putting for $A$, $B$, $C$ their values in terms of $\alpha$, $\beta$, $\gamma$ this is
\[
- 4\sin(\beta - \gamma)\sin(\gamma - \alpha)\sin(\alpha - \beta) = \sin\,(2\beta - 2\gamma)\,(1 - \cos 2\alpha)
\]
\[
+ \sin\,(2\gamma - 2\alpha)\,(1 - \cos 2\beta)
\]
\[
+ \sin\,(2\alpha - 2\beta)\,(1 - \cos 2\gamma),
\]

\newpage

%===============================================================
% PDF p.170 = book p.148
%===============================================================
\begin{center}
\small 148\hfill\textsc{on the area of the conic section \&c.}\hfill[192
\end{center}

\medskip

\noindent which is an identical equation; it is most readily proved by writing $x$, $y$, $z$ for $\tan\alpha$, $\tan\beta$, $\tan\gamma$; the equation thus becomes
\[
\frac{-4}{(1 + x^{2})(1 + y^{2})(1 + z^{2})}\,(y - z)(z - x)(x - y) = \Sigma\,\frac{1}{(1 + y^{2})(1 + z^{2})}\,[2 y(1 - z^{2}) - 2 z(1 - y^{2})]\,\frac{2 x^{2}}{1 + x^{2}},
\]
or multiplying out
\[
-(y - z)(z - x)(x - y) = \Sigma\,(y - z)(1 + y z)\,x^{2} = \Sigma x^{2}(y - z) + x y z\,\Sigma x\,(y - z),
\]
that is
\[
-(y - z)(z - x)(x - y) = \Sigma x^{2}(y - z) = x^{2}(y - z) + y^{2}(z - x) + z^{2}(x - y),
\]
which is an identity.

\medskip

\noindent [A different investigation of the formula $r^{2} = L\xi^{2} + M\eta^{2} + N\zeta^{2}$, by Dr Ferrers, was appended to the original Paper.]

\newpage

%===============================================================
% PDF p.171 = book p.149 -- Paper 193 begins
%===============================================================
\begin{center}
\small 193]\hfill\textsc{}\hfill\textsc{149}
\end{center}

\bigskip

\begin{center}
{\Large\bfseries 193.}\\[8pt]
{\large ON RODRIGUES' METHOD FOR THE ATTRACTION OF\\
ELLIPSOIDS.}\\[6pt]
\small [From the \emph{Quarterly Mathematical Journal}, vol.~\textsc{ii}.\ (1858), pp.~333--337.]
\end{center}

\bigskip

\noindent \textsc{The} following is in substance the method given in the ``M\'emoire sur l'attraction des Sph\'ero\"ides,'' par M.~Rodrigues, \emph{Corresp.\ sur l'\'Ecole Polyt.}, t.~\textsc{iii}.\ pp.~361--385 (1815). It will be seen that the method is very similar to that given two years before by Gauss, see my paper ``On Gauss' Method for the Attraction of Ellipsoids,'' \emph{Journal}, t.~\textsc{i}.\ pp.~162--166 [164]: the solution in fact depends upon the geometrical theorem therein quoted, viz.\ if $M$ be any point, $P$ a point of a closed surface, $PQ$ the normal (lying outside the surface) at the point $P$, $dS$ the element of the surface at that point, and if $MQ$ denotes the angle $MPQ$ and $\overline{MP}$ the distance of the points $M$ and $P$, then, theorem, the integral
\[
\iint \frac{dS\,\cos MQ}{\overline{MP^{2}}}
\]
has for its value
\[
0,\quad -2\pi \text{ or } -4\pi
\]
according as $M$ is exterior to, upon, or interior to the surface.

Suppose that $M$ is the attracted point and taking $A$, $B$, $C$ for the semiaxes of the surface of the attracting ellipsoid, or, if we please, for any semiaxes of an arbitrary ellipsoidal surface confocal with the surface of the attracting ellipsoid, let $P$ be a point on the surface of the interior similar ellipsoid whose semiaxes are $r A$, $r B$, $r C$. The coordinates of $M$ are taken to be $a$, $b$, $c$, and those of $P$ are taken to be $x$, $y$, $z$, and the value of the potential is
\[
V = \int\frac{d m}{\overline{MP}},
\]
where $d m$ is the element of mass.

\newpage

%===============================================================
% PDF p.172 = book p.150
%===============================================================
\begin{center}
\small 150\hfill\textsc{on rodrigues' method for the attraction of ellipsoids.}\hfill[193
\end{center}

\medskip

\noindent We may write
\[
x = r A\xi,\quad y = r B\eta,\quad z = r C\zeta,
\]
and then $\xi$, $\eta$, $\zeta$ will be the coordinates of a point $P'$ on the surface of a sphere, radius unity, corresponding in a definite manner to the point $P$ on the surface of the internal similar ellipsoid. And if $d\sigma$ be the element of the spherical surface, then we have
\[
d m = A B C\,r^{2}\,dr\,d\sigma,
\]
and therefore
\[
\frac{V}{A B C} = \int\frac{r^{2}\,dr\,d\sigma}{\overline{MP}};
\]
where, in order to obtain the value of the potential $V$ for the ellipsoid whose semiaxes are $A$, $B$, $C$, the integrations must be extended over the spherical surface and from $r = 0$ to $r = 1$.

Suppose that $dS$ is the element of the internal similar surface at $P$, and let $p$ be the perpendicular from the centre upon the tangent plane at $P$, we have
\[
dS = \frac{r^{2} A B C}{p}\,d\sigma.
\]
Let $P_{0}$ be the point on the ellipsoid $(A,\ B,\ rC)$ similarly situated to the point $P$ on the ellipsoid $(r A,\ r B,\ r C)$; the coordinates of $P_{0}$ are $A\xi$, $B\eta$, $C\zeta$; and if $p_{0}$ be the perpendicular from the centre upon the tangent plane at $P_{0}$, then $p = r p_{0}$, and the preceding equation becomes
\[
dS = \frac{r^{3} A B C}{p_{0}}\,d\sigma.
\]

Imagine now an ellipsoidal surface confocal with the surface $(A,\ B,\ C)$ and having for its semiaxes
\[
A + \delta A,\quad B + \delta B,\quad C + \delta C,
\]
and let $P_{0}'$ be the point on this surface which corresponds with the point $P_{0}$ on the surface $(A,\ B,\ C)$; that is, let $P_{0}'$ be the point whose coordinates are
\[
(A + \delta A)\,\xi,\quad (B + \delta B)\,\eta,\quad (C + \delta C)\,\zeta;
\]
and let $P'$ be in like manner the point whose coordinates are
\[
r\,(A + \delta A)\,\xi,\quad r\,(B + \delta B)\,\eta,\quad r\,(C + \delta C)\,\zeta;
\]

\newpage

%===============================================================
% PDF p.173 = book p.151
%===============================================================
\begin{center}
\small 193]\hfill\textsc{on rodrigues' method for the attraction of ellipsoids.}\hfill\textsc{151}
\end{center}

\medskip

\noindent the points $P$, $P'$ will be in like manner \emph{corresponding} points on the surface $(r A,\ r B,\ r C)$ and on the confocal surface $[r\,(A + \delta A),\ r\,(B + \delta B),\ r\,(C + \delta C)]$; and if the normal distance at the point $P_{0}$ of the first two surfaces is $\delta N$, then the normal distance at the point $P$ of the second two surfaces will be $r\,\delta N$. The decrement of $\overline{MP}$ will be equal to the normal distance $r\,\delta N$ of the two surfaces at the point $P$ multiplied into the cosine of the angle $MQ$, and we have, by a property of confocal surfaces,
\[
A\delta A = B\delta B = C\delta C = p_{0}\,\delta N = (\text{suppose})\,\tfrac{1}{2}\delta\theta,
\]
so we have therefore
\[
d\,\overline{MP} = -\,\frac{\tfrac{1}{2} r\,\delta\theta}{p_{0}}\,\cos MQ.
\]
Hence from the equation
\[
\frac{V}{A B C} = \int\frac{r^{2}\,dr\,d\sigma}{\overline{MP}}.
\]
We deduce
\[
\delta\frac{V}{A B C} = \int r^{2}\,dr\,d\sigma\,\frac{\tfrac{1}{2} r\,\delta\theta}{p_{0}}\,\frac{\cos MQ}{\overline{MP^{2}}}.
\]
But we have
\[
\frac{r^{2}\,d\sigma}{p_{0}} = \frac{dS}{A B C};
\]
and the equation thus becomes
\[
\delta\frac{V}{A B C} = \frac{\tfrac{1}{2}\,\delta\theta}{A B C}\int r\,dr\,\frac{dS\,\cos MQ}{\overline{MP^{2}}}.
\]

It may be proper to remark here by way of recapitulation that the course of the investigation has been as follows: viz.\ that, with a view to obtaining the potential $V$ of an attracting ellipsoid, we have found the increment of $\dfrac{V}{A B C}$ in passing from the ellipsoidal surface $(A,\ B,\ C)$ to the ellipsoidal surface $(A + \delta A,\ B + \delta B,\ C + \delta C)$, each of them confocal with the surface of the attracting ellipsoid; and that for finding such increment we have had to consider the two surfaces $(r A,\ r B,\ r C)$ and $[r\,(A + \delta A),\ r\,(B + \delta B),\ r\,(C + \delta C)]$ confocal to each other and respectively similar to the first-mentioned two surfaces.

Resuming the formula just obtained, the integral with respect to $dS$ is taken over the entire surface of the internal similar ellipsoid $(r A,\ r B,\ r C)$, and if the attracted point $M$ is external to the ellipsoid $(A,\ B,\ C)$ it will be external to the interior similar ellipsoid $(r A,\ r B,\ r C)$; hence in this case the double integral vanishes for all values of $r$, or we have
\[
\delta\frac{V}{A B C} = 0;
\]

\newpage

%===============================================================
% PDF p.174 = book p.152
%===============================================================
\begin{center}
\small 152\hfill\textsc{on rodrigues' method for the attraction of ellipsoids.}\hfill[193
\end{center}

\medskip

\noindent that is the function $V \div A B C$, which represents the ratio of the potential to the mass, is not altered in passing from the ellipsoid $(A,\ B,\ C)$ to the confocal ellipsoid
\[
(A + \delta A,\ B + \delta B,\ C + \delta C),
\]
or, what is the same thing, the potentials (and therefore the attractions) of confocal ellipsoids upon the same external point are proportional to their masses; this is in fact Maclaurin's theorem for the attraction of ellipsoids upon an external point.

But if the attracted point $M$ is interior to the ellipsoid $(A,\ B,\ C)$, then writing
\[
\frac{a^{2}}{A^{2}} + \frac{b^{2}}{B^{2}} + \frac{c^{2}}{C^{2}} = r'^{2};
\]
where $r'$ is less than unity, the double integral is $= 0$ from $r = 0$ to $r = r'$ and is $= -4\pi$ from $r = r'$ to $r = 1$, and we have
\begin{align*}
\delta\frac{V}{A B C} &= -\,\frac{2\pi\,\delta\theta}{A B C}\,\int_{r'}^{1}\,r\,dr \\
&= -\,\frac{\pi\,\delta\theta}{A B C}\,(1 - r'^{2}) \\
&= -\,\frac{\pi\,\delta\theta}{A B C}\left(\frac{a^{2}}{A^{2}} + \frac{b^{2}}{B^{2}} + \frac{c^{2}}{C^{2}} - 1\right);
\end{align*}
that is, the right-hand side of the equation is the increment (or taken with its sign reversed so as to be positive, it is the decrement) of the function $V \div A B C$ in passing from the ellipsoid $(A,\ B,\ C)$ to the confocal ellipsoid
\[
(A + \delta A,\ B + \delta B,\ C + \delta C),
\]
where
\[
\tfrac{1}{2}\delta\theta = A\delta A = B\delta B = C\delta C.
\]

The preceding formula gives at once the potential for an interior point; in fact taking $\alpha$, $\beta$, $\gamma$ for the semiaxes of the ellipsoid and writing
\[
A^{2} = \alpha^{2} + \theta,\quad B^{2} = \beta^{2} + \theta,\quad C^{2} = \gamma^{2} + \theta,
\]
and using the ordinary symbol $d$ instead of $\delta$, we have
\[
\frac{d}{d\theta}\,\frac{V}{\sqrt{(\alpha^{2} + \theta)(\beta^{2} + \theta)(\gamma^{2} + \theta)}} = \frac{\pi}{\sqrt{(\alpha^{2} + \theta)(\beta^{2} + \theta)(\gamma^{2} + \theta)}}\left\{\frac{a^{2}}{\alpha^{2} + \theta} + \frac{b^{2}}{\beta^{2} + \theta} + \frac{c^{2}}{\gamma^{2} + \theta} - 1\right\},
\]
and integrating from $\theta = 0$ to $\theta = \infty$, we have
\[
-\,\frac{V}{\alpha\beta\gamma} = \pi\,\int_{0}^{\infty}\,\frac{d\theta}{\sqrt{(\alpha^{2} + \theta)(\beta^{2} + \theta)(\gamma^{2} + \theta)}}\left\{\frac{a^{2}}{\alpha^{2} + \theta} + \frac{b^{2}}{\beta^{2} + \theta} + \frac{c^{2}}{\gamma^{2} + \theta} - 1\right\},
\]

\newpage

%===============================================================
% PDF p.175 = book p.153
%===============================================================
\begin{center}
\small 193]\hfill\textsc{on rodrigues' method for the attraction of ellipsoids.}\hfill\textsc{153}
\end{center}

\medskip

\noindent where $V$ is now the potential for the ellipsoid whose semiaxes are $\alpha$, $\beta$, $\gamma$; and we have therefore
\[
V = -\,\pi\alpha\beta\gamma\,\int_{0}^{\infty}\,\frac{d\theta}{\sqrt{(\alpha^{2} + \theta)(\beta^{2} + \theta)(\gamma^{2} + \theta)}}\left\{\frac{a^{2}}{\alpha^{2} + \theta} + \frac{b^{2}}{\beta^{2} + \theta} + \frac{c^{2}}{\gamma^{2} + \theta} - 1\right\}.
\]
To find the potential for an external point it is only necessary to remark that by the theorem above demonstrated $V \div \alpha\beta\gamma$ is equal to the corresponding function for the confocal ellipsoid through the attracted point, that is for the ellipsoid whose semiaxes are $\sqrt{(\alpha^{2} + \theta_{1})}$, $\sqrt{(\beta^{2} + \theta_{1})}$, $\sqrt{(\gamma^{2} + \theta_{1})}$, where $\theta_{1}$ is a positive quantity such that
\[
\frac{a^{2}}{\alpha^{2} + \theta_{1}} + \frac{b^{2}}{\beta^{2} + \theta_{1}} + \frac{c^{2}}{\gamma^{2} + \theta_{1}} = 1;
\]
hence in the value of $V \div \alpha\beta\gamma$ we have only to write the above values in the place of $\alpha$, $\beta$, $\gamma$; and if we then write $\theta - \theta_{1}$ in the place of $\theta$ the limits will be $\infty$, $\theta_{1}$, and the expression for the potential is
\[
V = -\,\pi\alpha\beta\gamma\,\int_{\theta_{1}}^{\infty}\,\frac{d\theta}{\sqrt{(\alpha^{2} + \theta)(\beta^{2} + \theta)(\gamma^{2} + \theta)}}\left\{\frac{a^{2}}{\alpha^{2} + \theta} + \frac{b^{2}}{\beta^{2} + \theta} + \frac{c^{2}}{\gamma^{2} + \theta} - 1\right\};
\]
this completes the investigation.

\end{document}
